\documentclass[aspectratio=169,10pt,t]{beamer}
\usetheme{faeder}

\usepackage{hyperref}
\usepackage{qrcode}
\hypersetup{colorlinks=true,urlcolor=fsTeal,linkcolor=fsNavy}

\graphicspath{{figures/}}
\newcommand{\courserepo}{https://github.com/jrfaeder/RBM_Intro_BioNetGen}
\newcommand{\lecturerepo}{https://github.com/jrfaeder/SystemsImmunology}
\newcommand{\bngplayground}{https://ruleworld.github.io/bngplayground/}
\newcommand{\bookurl}{https://ruleworld.github.io/bionetgen-book/}
\newcommand{\yes}{{\color{fsGreen}\checkmark}}
\newcommand{\no}{{\color{fsAccent}$\times$}}

% ---------------------------------------------------------------- figures
% Figures are pages of the 2024 PowerPoint deck (figures/slideNN.pdf, 960 x 540 pt).
% Regions are given in slide points, (x1,y1) top left to (x2,y2) bottom right,
% as read from the .pptx picture positions.
%
%   \regionfig[opts]{NN}{x1}{y1}{x2}{y2}            crop of slide NN
%   \maskedfig[opts]{NN}{x1}{y1}{x2}{y2}{masks}     same, with masks drawn on top
%   \hide{spec}{x1}{y1}{x2}{y2}   white mask on the slides in spec, never in the handout;
%                                 use inside \maskedfig to reveal a figure one panel at a time
%   \erase{x1}{y1}{x2}{y2}        white mask everywhere (removes old callouts)
\newcommand{\regionfig}[6][]{%
  \pgfmathsetmacro\rfB{540-#6}\pgfmathsetmacro\rfR{960-#5}%
  \includegraphics[trim={#3bp \rfB bp \rfR bp #4bp},clip,#1]{slide#2}}
\newcommand{\maskedfig}[7][]{%
  \begin{tikzpicture}
    \node[anchor=south west,inner sep=0] (img) at (0,0) {\regionfig[#1]{#2}{#3}{#4}{#5}{#6}};
    \begin{scope}[x={(img.south east)},y={(img.north west)}]
      \def\cropL{#3}\def\cropT{#4}\def\cropR{#5}\def\cropB{#6}
      #7
    \end{scope}
  \end{tikzpicture}}
\newcommand{\maskcoords}[4]{%
  \pgfmathsetmacro\mxa{(#1-\cropL)/(\cropR-\cropL)}%
  \pgfmathsetmacro\mxb{(#3-\cropL)/(\cropR-\cropL)}%
  \pgfmathsetmacro\mya{1-(#2-\cropT)/(\cropB-\cropT)}%
  \pgfmathsetmacro\myb{1-(#4-\cropT)/(\cropB-\cropT)}}
\newcommand{\hide}[5]{\maskcoords{#2}{#3}{#4}{#5}%
  \fill[white,visible on=<#1|handout:0>] (\mxa,\mya) rectangle (\mxb,\myb);}
\newcommand{\erase}[4]{\maskcoords{#1}{#2}{#3}{#4}%
  \fill[white] (\mxa,\mya) rectangle (\mxb,\myb);}
% \replacefig{file}{x1}{y1}{x2}{y2}: inside \maskedfig, cover that region of the old slide
% and draw a corrected image in its place, scaled to fit (aspect ratio kept). Put it
% before any \hide masks so the masks still cover it on earlier steps.
\newcommand{\replacefig}[5]{\maskcoords{#2}{#3}{#4}{#5}%
  \fill[white] (\mxa,\mya) rectangle (\mxb,\myb);
  \pgfpointanchor{img}{north east}\pgfgetlastxy{\imgW}{\imgH}%
  \pgfmathsetlengthmacro\rpW{(\mxb-\mxa)*\imgW}%
  \pgfmathsetlengthmacro\rpH{(\mya-\myb)*\imgH}%
  \node[inner sep=0] at ({(\mxa+\mxb)/2},{(\mya+\myb)/2})
    {\includegraphics[width=\rpW,height=\rpH,keepaspectratio]{#1}};}
% Same, but only if figures/<name>.pdf or figures/<name>.png exists; otherwise the
% original slide content is left alone.
\newcommand{\replacefigifexists}[5]{%
  \IfFileExists{figures/#1.pdf}{\replacefig{#1.pdf}{#2}{#3}{#4}{#5}}{%
    \IfFileExists{figures/#1.png}{\replacefig{#1.png}{#2}{#3}{#4}{#5}}{}}}

% Discussion prompt box
\newcommand{\discuss}[1]{%
  \begin{tikzpicture}
    \node[fill=fsGold!18,draw=fsGold,rounded corners=2pt,inner sep=2.2mm,
          text width=\dimexpr\linewidth-4.4mm\relax,font=\small] {\textbf{Discuss.} #1};
  \end{tikzpicture}}

\title[Mechanistic Modeling of Intracellular Signaling]{Mechanistic Modeling of\\Intracellular Signaling}
\subtitle{Advances in Systems Immunology}
\author{Jim Faeder}
\institute{Department of Computational and Systems Biology\\University of Pittsburgh}
\date{October 14, 2026}

\begin{document}

% ======================================================================
\begin{frame}[plain]
  \titlepage
  % QR code to the lecture repository, lower right of the title page
  \begin{tikzpicture}[remember picture,overlay]
    \node[anchor=south east,fill=white,rounded corners=2pt,inner sep=1.5mm] (qr)
      at ([xshift=-12mm,yshift=11mm]current page.south east)
      {\qrcode[height=1.9cm,nolink]{\lecturerepo}};
    \node[anchor=north,font=\tiny,text=white!80!fsNavy,inner sep=0] at ([yshift=-1.2mm]qr.south)
      {github.com/jrfaeder/SystemsImmunology};
  \end{tikzpicture}
\end{frame}

\begin{frame}{Goals for today}
  By the end of this lecture you should be able to
  \begin{enumerate}[<+->]
    \setlength\itemsep{1.5mm}
    \item explain what a \hl{mechanistic model} is and what it is good for
    \item translate a reaction scheme into \hl{rate equations}
    \item explain why \hl{combinatorial complexity} motivates rule-based modeling
    \item describe how models are \hl{calibrated} to data and used to make predictions
    \item follow a \hl{model-selection} argument that identifies the minimal signaling
          architecture consistent with T cell data
    \item judge how far mechanistic models currently \hl{scale}
  \end{enumerate}
\end{frame}

% ======================================================================
% Part 1  (~10 min)
\section{Why build mechanistic models?}

\begin{frame}{Systems biology closes the loop between experiment and model}
  \begin{columns}[c]
    \begin{column}{0.46\textwidth}
      \centering
      \maskedfig[height=0.66\textheight]{05}{300}{190}{640}{527}{%
        \erase{300}{190}{368}{222}\erase{368}{190}{375}{212}\erase{300}{190}{640}{196}% old slide text
        \erase{626}{285}{640}{405}}% edge of the old "Key point" box
    \end{column}
    \begin{column}{0.5\textwidth}
      \small
      \textbf{Goals}: \emph{map} the networks that control cell behavior,
      \emph{explain} how they make decisions, and \emph{manipulate} or
      \emph{engineer} them to treat disease.
      \begin{itemize}
        \setlength\itemsep{1.5mm}
        \item<2-> ``Hypothesis-driven research'': a closed loop of \emph{wet} and
              \emph{dry} experiments.
        \item<3-> Models turn current knowledge into predictions; experiments test them;
              the results revise the model.
      \end{itemize}
      \vskip2mm
      \onslide<4->
      \begin{block}{Key point}
        Computational models of biological systems are needed to close the loop.
      \end{block}
    \end{column}
  \end{columns}
  \source{H. Kitano, \emph{Science} \textbf{295}, 1662 (2002)}
\end{frame}

\begin{frame}{A diagram is a model; a computational model formalizes it}
  \centering
  % Left: slide 8 of the 2024 deck (diagram, then the experimental biologist pondering it).
  \maskedfig[height=0.5\textheight]{08}{144}{69}{947}{516}{%
    \erase{144}{69}{382}{110}% old slide title
    \hide{1}{330}{60}{950}{520}}%
  \hspace{2mm}%
  % Right: slide 9 from the first arrow on (computational biologist, BioNetGen model,
  % then the corrected dose-response plot).
  \maskedfig[height=0.5\textheight]{09}{309}{105}{840}{540}{%
    % corrected d-r plot; box starts below the arrow tip (y=343) so the arrow is not covered,
    % and the strip of the old plot beside the arrow (y 335-343) is erased separately
    \erase{608}{335}{714}{343}\erase{752}{335}{840}{343}%
    \replacefig{slide06-dose-response.pdf}{608}{344}{840}{540}%
    \hide{1-2}{305}{100}{845}{545}%
    \hide{1-3}{600}{300}{845}{545}}
  \vskip2mm
  \raggedright\small
  \begin{itemize}
    \setlength\itemsep{1mm}
    \item<2-> From a pathway diagram, an immunologist reasons out experiments and expected
          results: informal modeling.
    \item<4-> A computational model makes the same reasoning precise, and the computer works
          out the consequences, past the limits of intuition.
  \end{itemize}
\end{frame}

\begin{frame}{What good is a model?}
  \begin{columns}[T]
    \begin{column}{0.58\textwidth}
      \begin{itemize}
        \setlength\itemsep{3mm}
        \item<1-> ``If the use of a model stimulates experiments, even if the model is
              subsequently shown to be wrong, then it has been successful.''\\
              {\small\color{fsGray}--- J. D. Murray}
        \item<2-> ``All models are wrong, but some are useful.''\\
              {\small\color{fsGray}--- George E. P. Box}
        \item<3-> ``What I cannot create, I do not understand.''\\
              {\small\color{fsGray}--- Richard Feynman}
      \end{itemize}
    \end{column}
    \begin{column}{0.38\textwidth}
      \small
      \begin{itemize}
        \setlength\itemsep{3mm}
        \item<4-> The main goal of modeling is to understand the system better, and to use
              that understanding to design new experiments.
        \item<5-> Every model simplifies, so every model is wrong somewhere. What matters is
              whether it is useful for the question you are asking.
        \item<6-> Building a model tests our understanding: if we cannot build a model that
              works, we do not yet know how the system works.
      \end{itemize}
    \end{column}
  \end{columns}
\end{frame}

\begin{frame}{Reverse engineering biological systems is hard}
  \begin{columns}[c]
    \begin{column}{0.4\textwidth}
      \centering
      \regionfig[width=\linewidth]{12}{187}{132}{771}{539}\\
      {\scriptsize EGFR (ErbB1--4) signaling map}
    \end{column}
    \begin{column}{0.24\textwidth}
      \centering
      \onslide<2->
      \regionfig[width=\linewidth]{13}{134}{118}{430}{495}
    \end{column}
    \begin{column}{0.32\textwidth}
      \small
      \onslide<2->{Can a biologist fix a radio? Compared with an engineer:}
      \begin{itemize}
        \setlength\itemsep{1mm}
        \item<3-> We don't know all the parts.
        \item<4-> We don't know how they are connected.
        \item<5-> Many parts belong to several circuits.
        \item<6-> We don't know the design principles.
        \item<7-> We lack a common language.
      \end{itemize}
    \end{column}
  \end{columns}
  \source{Y. Lazebnik, \emph{Cancer Cell} \textbf{2}, 179 (2002)}
\end{frame}

\begin{frame}{Even simple systems can confuse us}
  \begin{columns}[c]
    \begin{column}{0.64\textwidth}
      \centering
      \maskedfig[width=\linewidth]{14}{144}{106}{806}{532}{%
        \hide{1}{460}{140}{810}{535}}
    \end{column}
    \begin{column}{0.33\textwidth}
      \small
      \begin{itemize}
        \setlength\itemsep{2mm}
        \item<1-> A three-step cascade relaxes smoothly to steady state.
        \item<2-> Add one negative feedback loop, and the same components can
              \hl{oscillate}.
        \item<3-> Oscillation is \hl{emergent}: no component oscillates by itself. It
              comes from the interactions, once the parts are assembled into a network.
      \end{itemize}
    \end{column}
  \end{columns}
\end{frame}

% ======================================================================
% Part 2  (~15 min)
\section{From mechanisms to equations}

\begin{frame}{The standard modeling protocol}
  \begin{enumerate}
    \setlength\itemsep{1.5mm}
    \item<1-> Identify the components and their interactions.
          \vskip1mm
          {\centering \regionfig[height=0.26\textheight]{16}{204}{144}{780}{276}\par}
    \item<2-> Write the reactions, and from them the rate equations:
          \quad $\dot{\mathbf{x}} = \mathbf{S}\cdot\mathbf{v}(\mathbf{x})$
    \item<3-> Determine the concentrations and rate constants.
    \item<4-> Simulate and analyze the model: ODEs, stochastic simulation (SSA),
          PDEs for space, Brownian dynamics, \ldots
  \end{enumerate}
\end{frame}

\begin{frame}{From reactions to differential equations}
  \begin{columns}[T]
    \begin{column}{0.6\textwidth}
      Ligand $L$ binds receptor $R$:
      \[ R + L \xrightarrow{\;k_1\;} RL \]
      \onslide<2->
      Mass action: the rate is proportional to the reactant concentrations,
      \[ v_1 = k_1\,[R]\,[L] \]
      \onslide<3->
      Each species changes by the fluxes that make it minus those that consume it:
      \[
        \frac{d[R]}{dt} = -v_1 + \cdots \quad
        \frac{d[L]}{dt} = -v_1 + \cdots \quad
        \frac{d[RL]}{dt} = +v_1 + \cdots
      \]
    \end{column}
    \begin{column}{0.36\textwidth}
      \onslide<1->
      \centering
      \regionfig[height=0.34\textheight]{17}{646}{89}{804}{268}
      \vskip2mm
      \raggedright\small
      \onslide<4->
      The ``$\cdots$'' are terms from other reactions. Fluxes combine by \hl{addition}.
      \vskip2mm
      \onslide<5->
      In general: $\dot{\mathbf{x}} = \mathbf{S}\cdot\mathbf{v}(\mathbf{x})$, with
      stoichiometry matrix $\mathbf{S}$ and rate vector $\mathbf{v}$.
    \end{column}
  \end{columns}
\end{frame}

\begin{frame}{The Kholodenko model explains transient EGFR phosphorylation}
  \centering
  \regionfig[height=0.46\textheight]{18}{139}{124}{422}{509}%
  \hspace{3mm}%
  \maskedfig[height=0.46\textheight]{19}{43}{135}{948}{534}{%
    \hide{1}{40}{130}{952}{538}%
    \hide{1-2}{600}{130}{952}{538}}
  \vskip2mm
  \raggedright\small
  \begin{itemize}
    \setlength\itemsep{1mm}
    \item<1-> A reaction network written as one ODE per species: 22 species, 25 reactions.
    \item<2-> Why does receptor phosphorylation fall so quickly after EGF stimulation?
    \item<3-> Model: phosphorylated Shc has lower affinity for pEGFR.
  \end{itemize}
  \source{B. N. Kholodenko et al., \emph{J.\ Biol.\ Chem.}\ \textbf{274}, 30169 (1999)}
\end{frame}

\begin{frame}{The Hoffmann model of NF-$\kappa$B: negative feedback shapes the response}
  \centering
  \maskedfig[height=0.6\textheight]{20}{41}{109}{949}{539}{%
    \hide{1}{480}{105}{955}{358}%
    \hide{1-2}{480}{358}{955}{540}}
  \only<3>{}% last mask lifts on slide 3
  \source{A. Hoffmann et al., \emph{Science} \textbf{298}, 1241 (2002)}
\end{frame}

\begin{frame}{Kinetic proofreading: how T cells tell ligands apart}
  \begin{columns}[T]
    \begin{column}{0.5\textwidth}
      \centering
      \regionfig[width=0.9\linewidth]{21}{249}{133}{711}{249}
      \vskip2mm
      \regionfig[width=0.75\linewidth]{21}{67}{291}{379}{492}
    \end{column}
    \begin{column}{0.47\textwidth}
      \small
      \begin{itemize}
        \setlength\itemsep{1.5mm}
        \item<1-> The bound TCR must pass through $N$ modification steps, rate $k_p$,
              before it signals. If the ligand unbinds, the receptor resets.
        \item<2-> Fraction of binding events that reach the signaling state:
              \[ \left(\frac{k_p}{k_p + k_{\mathrm{off}}}\right)^{N} \approx
                 e^{-k_{\mathrm{off}}\,\tau}, \qquad \tau = N/k_p \]
        \item<3-> \hl{Key idea}: the delay $\tau$ allows error correction.
        \item<4-> A delay of 3 s amplifies a 10-fold difference in off-rate
              ($k_{\mathrm{off}}=0.5$ vs.\ 5 s$^{-1}$) into a $\sim10^{6}$-fold difference
              in signal.
        \item<5-> Fc$\varepsilon$RI, BCR, and TCR signal through sequential phosphorylation
              and recruitment steps, each a candidate proofreading stage.
      \end{itemize}
    \end{column}
  \end{columns}
  \source{T. W. McKeithan, \emph{PNAS} \textbf{92}, 5042 (1995); Goldstein et al., \emph{Nat.\ Rev.\ Immunol.}\ \textbf{4}, 445 (2004)}
\end{frame}

\begin{frame}{Detailed models of immunoreceptor signaling become complex}
  \centering
  \maskedfig[height=0.56\textheight]{24}{17}{137}{948}{521}{%
    \hide{1}{540}{130}{952}{525}}
  \vskip1mm
  \raggedright\small
  \onslide<1->{Left: TCR discrimination with SHP-1 and ERK feedback \emph{(network complexity)}.}
  \onslide<2->{Right: Fc$\varepsilon$RI, where a few interactions produce hundreds of complexes
  \emph{(combinatorial complexity)}.}
  \source{Altan-Bonnet and Germain, \emph{PLoS Biol.}\ (2005); Faeder et al., \emph{J.\ Immunol.}\ (2003)}
\end{frame}

% ======================================================================
% Part 3  (~12 min)
\section{Combinatorial complexity and rule-based modeling}

\begin{frame}{Combinatorial complexity}
  \centering
  \maskedfig[height=0.6\textheight]{25}{114}{91}{814}{540}{%
    \hide{1}{110}{248}{820}{545}%
    \hide{1-2}{550}{248}{820}{545}}
  \vskip1mm
  \raggedright\small
  \onslide<3->{Four molecules, a handful of interactions: \hl{354 species, 3,680 reactions}.}
  \source{Faeder et al., \emph{J.\ Immunol.}\ \textbf{170}, 3769 (2003)}
\end{frame}

\begin{frame}[fragile]{Rule-based modeling: describe the interactions, generate the equations}
  \begin{columns}[T]
    \begin{column}{0.36\textwidth}
      \centering
      \regionfig[width=0.85\linewidth]{29}{180}{144}{381}{343}
      \vskip2mm
      \raggedright\small
      \onslide<2->
      Molecules are objects with \textbf{components} (binding and modification sites),
      which can have \textbf{states}.
    \end{column}
    \begin{column}{0.6\textwidth}
      \onslide<2->
\begin{lstlisting}[language=bngl,basicstyle=\ttfamily\scriptsize]
begin molecule types
  Lig(l,l)
  Lyn(U,SH2)
  Syk(tSH2,l~U~P,a~U~P)
  Rec(a,b~U~P,g~U~P)
end molecule types
\end{lstlisting}
      \onslide<3->
\begin{lstlisting}[language=bngl,basicstyle=\ttfamily\scriptsize]
begin reaction rules
  # Ligand-receptor binding
  Rec(a) + Lig(l,l) <-> Rec(a!1).Lig(l!1,l)  kp1, km1
  # Receptor aggregation
  Rec(a) + Lig(l,l!1) <-> Rec(a!2).Lig(l!2,l!1)  kp2, km2
  # Constitutive Lyn-receptor binding
  Rec(b~U) + Lyn(U,SH2) <-> Rec(b~U!1).Lyn(U!1,SH2)  kpL, kmL
  ...
end reaction rules
\end{lstlisting}
      \small
      \onslide<4->
      A rule acts on every species containing its pattern (``don't write, don't care'').
    \end{column}
  \end{columns}
\end{frame}

\begin{frame}{Network generation}
  \begin{columns}[c]
    \begin{column}{0.5\textwidth}
      \centering
      \begin{tikzpicture}[box/.style={draw=fsNavy,fill=fsLight,rounded corners=2pt,
          minimum width=2.6cm,minimum height=8mm,align=center,font=\small},
          arr/.style={-{Stealth[length=2mm]},line width=1pt,fsGray}]
        \node[box] (m) at (0,0) {Molecule types\\\textbf{4}};
        \node[box] (r) at (3.2,0) {Reaction rules\\\textbf{19}};
        \node[box,fill=fsGold!25,draw=fsGold,visible on=<2->] (g) at (1.6,-1.6) {Network generation};
        \node[box,visible on=<3->] (s) at (0,-3.2) {Species\\\textbf{354}};
        \node[box,visible on=<3->] (x) at (3.2,-3.2) {Reactions\\\textbf{3,680}};
        \draw[arr,visible on=<2->] (m) -- (g); \draw[arr,visible on=<2->] (r) -- (g);
        \draw[arr,visible on=<3->] (g) -- (s); \draw[arr,visible on=<3->] (g) -- (x);
      \end{tikzpicture}
    \end{column}
    \begin{column}{0.46\textwidth}
      \begin{itemize}
        \setlength\itemsep{2mm}
        \item<2-> Apply each rule to the seed species, add the new species and reactions,
              and repeat until nothing new appears.
        \item<3-> The Fc$\varepsilon$RI model: 19 rules generate 3,680 reactions.
        \item<4-> The generated network is simulated with ODEs or SSA, as before.
      \end{itemize}
    \end{column}
  \end{columns}
  \source{Blinov et al., \emph{Bioinformatics} \textbf{20}, 3289 (2004); Harris et al., \emph{Bioinformatics} \textbf{32}, 3366 (2016)}
\end{frame}

\begin{frame}{Rules capture knowledge at large scale}
  \begin{columns}[c]
    \begin{column}{0.6\textwidth}
      \centering
      \maskedfig[width=\linewidth]{33}{33}{130}{652}{510}{%
        \hide{1}{300}{125}{656}{515}}
    \end{column}
    \begin{column}{0.36\textwidth}
      \small
      \begin{itemize}
        \setlength\itemsep{1.5mm}
        \item<1-> Rules are a precise encoding of the modeled structures and interactions.
        \item<2-> TCR model: 20 molecule types, 129 rules. The number of possible species is
              too large to generate the network at all.
        \item<3-> Rule-based models can be visualized and extended as knowledge grows.
      \end{itemize}
    \end{column}
  \end{columns}
  \source{Faeder et al., \emph{J.\ Immunol.}\ (2003); Chylek et al., \emph{PLOS One} (2014)}
\end{frame}

\begin{frame}{Network-free simulation}
  \begin{columns}[c]
    \begin{column}{0.56\textwidth}
      \centering
      \maskedfig[width=\linewidth]{34}{198}{132}{958}{524}{%
        \hide{1}{575}{132}{960}{526}}
    \end{column}
    \begin{column}{0.4\textwidth}
      \small
      \begin{itemize}
        \setlength\itemsep{2mm}
        \item<1-> \textbf{Indirect}: generate the network, then simulate it.
        \item<2-> \textbf{Direct} (NFsim): track individual molecules and apply rules
              as the simulation runs. Only complexes that actually form are represented.
        \item<3-> Cost scales with the number of rules and molecules, not with the size
              of the network.
      \end{itemize}
    \end{column}
  \end{columns}
  \source{Sneddon, Faeder, and Emonet, \emph{Nat.\ Methods} \textbf{8}, 177 (2011)}
\end{frame}

\begin{frame}{Tools: BioNetGen today}
  \begin{columns}[T]
    \begin{column}{0.64\textwidth}
      \begin{itemize}
        \setlength\itemsep{1.8mm}
        \item \textbf{BioNetGen}: \href{https://bionetgen.org}{bionetgen.org}, with a
              \textbf{VS Code extension} for writing and running models
        \item \textbf{BNG Playground}: write and simulate models in the browser,
              no installation\\
              \href{\bngplayground}{ruleworld.github.io/bngplayground}
        \item \textbf{bngsim}: ODE and stochastic simulation from Python
              (\code{pip install bngsim bionetgen})
        \item \textbf{Open book and course}: \emph{Rule-Based Modeling of Biological
              Systems}, with models, notebooks, and slides\\
              \href{\bookurl}{ruleworld.github.io/bionetgen-book}
      \end{itemize}
    \end{column}
    \begin{column}{0.32\textwidth}
      \centering
      \qrcode[height=2.3cm]{\bngplayground}\\[1mm]
      {\small BNG Playground}\\[3mm]
      \qrcode[height=2.3cm]{\bookurl}\\[1mm]
      {\small The book}
    \end{column}
  \end{columns}
\end{frame}

% ======================================================================
% Part 4  (~10 min)
\section{Calibrating models to data: T cell fate}

\begin{frame}{The systems modeling workflow}
  \centering
  % Reveal the cycle one stage at a time: 1 construct (box + contact map),
  % 2 calibrate, 3 analyze, 4 experiments, 5 arrow back to the model + "Repeat as needed"
  \maskedfig[height=0.78\textheight]{39}{129}{99}{877}{537}{%
    \hide{1}{560}{128}{880}{420}%                  arrow, Calibrate box, dose fits
    \hide{1-2}{557}{407}{639}{482}%                 arrow to Analyze
    \hide{1-2}{622}{453}{826}{541}%                 Analyze box
    \hide{1-2}{390}{377}{547}{540}%                 PTEN dose curve
    \hide{1-3}{302}{358}{386}{432}%                 arrow to experiments
    \hide{1-3}{126}{256}{314}{443}%                 Nur77 / pS6 data
    \hide{1-3}{126}{445}{340}{541}%                 Perform experiments box
    \hide{1-4}{257}{194}{340}{269}%                 arrow back to the model
    \hide{1-4}{370}{318}{568}{353}}%                "Repeat as needed"
  \only<5>{}% last masks lift on slide 5
\end{frame}

\begin{frame}{Parameter estimation: inferring model parameters from data}
  \centering
  \maskedfig[height=0.6\textheight]{40}{47}{133}{959}{502}{%
    \hide{1}{300}{125}{962}{505}%
    \hide{1-2}{660}{125}{962}{505}}
  \vskip1mm
  \raggedright\small
  \onslide<2->{Choose parameters $\theta$ to minimize a cost: the mismatch between
  simulated and measured data.}
  \onslide<3->{Most rate constants and concentrations have not been measured, so they are
  inferred this way.}
\end{frame}

\begin{frame}{T cell differentiation: an experiment--model collaboration}
  \centering
  \maskedfig[height=0.58\textheight]{41}{11}{104}{886}{533}{%
    \hide{1}{265}{100}{890}{536}%
    \hide{1-2}{690}{100}{890}{536}}
  \vskip1mm
  \raggedright\small
  \onslide<1->{Low antigen dose favors induction of Foxp3$^+$ regulatory T cells.}
  \onslide<3->{What mechanism turns dose into a fate decision?}
  \source{Miskov-Zivanov et al., \emph{Sci.\ Signal.}\ \textbf{6}, ra97 (2013)}
\end{frame}

\begin{frame}{A rule-based model of TCR--Akt signaling}
  \begin{columns}[c]
    \begin{column}{0.32\textwidth}
      \begin{itemize}
        \setlength\itemsep{2mm}
        \item<1-> 110 rules
        \item<2-> 125 species, 367 reactions
        \item<3-> 144 parameters, 81 of them fitted
      \end{itemize}
      \vskip3mm
      \onslide<4->
      {\small Fitted with Bayesian parameter estimation (parallel tempering MCMC), which
      gives an \emph{ensemble} of parameter sets consistent with the data.}
    \end{column}
    \begin{column}{0.66\textwidth}
      \onslide<1->
      \centering
      \regionfig[height=0.78\textheight]{42}{357}{115}{884}{540}
    \end{column}
  \end{columns}
\end{frame}

\begin{frame}{The model predicts sharp thresholds in antigen dose}
  \centering
  \regionfig[width=0.9\linewidth]{45}{8}{142}{945}{530}
  \vskip1mm
  \raggedright\small
  \begin{itemize}
    \item<2-> Akt and FoxO1 switch sharply with dose; PTEN shows a narrow window.
    \item<3-> An Akt--FoxO1--PTEN \hl{positive feedback loop} gives bistability.
  \end{itemize}
\end{frame}

\begin{frame}{Prediction: lowering PTEN lowers the threshold}
  \centering
  \regionfig[height=0.5\textheight]{47}{64}{121}{881}{489}
  \vskip1mm
  \raggedright\small
  \onslide<2->{Model: halving PTEN shifts the switch to lower dose and shorter stimulation.}\\
  \onslide<3->{Test: compare PTEN$^{+/+}$ and PTEN$^{+/-}$ mouse T cells.}
\end{frame}

\begin{frame}{Test: commitment time is reduced in PTEN$^{+/-}$ T cells}
  \centering
  \regionfig[height=0.58\textheight]{48}{84}{130}{849}{537}
  \vskip1mm
  \raggedright\small
  \onslide<2->{\hl{Strength and duration} of TCR signaling control commitment; PTEN level sets
  the threshold.}
\end{frame}

% ======================================================================
% Part 5  (~22 min)
\section{\texorpdfstring{Case study:\\T cell ligand discrimination}{Case study: T cell ligand discrimination}}

\begin{frame}{Case study: a minimal model from comprehensive data}
  \centering
  \regionfig[width=0.72\linewidth]{50}{148}{153}{812}{349}
  \vskip4mm
  \raggedright
  \begin{itemize}
    \setlength\itemsep{1.5mm}
    \item<2-> T cells respond to pMHC ligands that vary over a million-fold in affinity and in
          dose. What signaling architecture explains the full dose--response surface?
    \item<3-> Approach: \hl{phenotypic models}. Abstract signaling species fitted to data,
          not molecular detail, and rejected one by one.
  \end{itemize}
  \source{Lever et al., \emph{PNAS} \textbf{113}, E6630 (2016); Lever et al., \emph{Nat.\ Rev.\ Immunol.}\ \textbf{14}, 619 (2014)}
\end{frame}

\begin{frame}{Experimental system and cytokine response data}
  \centering
  \maskedfig[height=0.6\textheight]{51}{77}{104}{875}{540}{%
    \hide{1}{430}{120}{880}{305}%
    \hide{1-2}{70}{302}{880}{545}}
  \vskip1mm
  \raggedright\footnotesize
  \onslide<2->{Peptide variants span $K_D$ from $\sim$10$^{-10}$ to $>$10$^{-4}$ M.}
  \onslide<3->{Readouts: IFN-$\gamma$, MIP-1$\beta$ (primary CD8$^+$), IL-8 (Jurkat).}
\end{frame}

\begin{frame}{Key features of the data}
  \begin{enumerate}
    \setlength\itemsep{1mm}
    \small
    \item<1-> \hl{Bell-shaped} dose--response for high- but not low-affinity ligands.
    \item<2-> The peak amplitude of the bell is \hl{independent of affinity}.
    \item<3-> At low doses (left of the peak), a \emph{single} intermediate-affinity ligand
          gives the largest response.
    \item<4-> At high doses (right of the peak), \emph{different} intermediate-affinity
          ligands give the largest response.
  \end{enumerate}
  \vskip1mm
  \centering
  \regionfig[height=0.44\textheight]{52}{75}{290}{873}{512}
  \source{Lever et al., \emph{PNAS} (2016)}
\end{frame}

\begin{frame}{Sequential model rejection}
  \centering
  \maskedfig[width=0.6\linewidth]{56}{66}{144}{862}{473}{%
    \hide{1}{248}{140}{866}{476}%
    \hide{1-2}{436}{140}{866}{476}%
    \hide{1-3}{630}{140}{866}{476}}
  \vskip1mm
  \raggedright\footnotesize
  \renewcommand{\arraystretch}{1.1}%
  \begin{tabular}{@{}p{3mm}p{0.7\linewidth}p{0.2\linewidth}@{}}
    \uncover<1->{\textbf{a}} & \uncover<1->{Occupancy: sigmoidal, no bell}
      & \uncover<1->{\no{} feature 1} \\
    \uncover<2->{\textbf{b}} & \uncover<2->{+ incoherent feedforward (IFF): bell, but affinity
      only shifts the curve} & \uncover<2->{\no{} features 3, 4} \\
    \uncover<3->{\textbf{c}} & \uncover<3->{Kinetic proofreading + IFF: peak falls for the
      highest affinities} & \uncover<3->{\no{} feature 2} \\
    \uncover<4->{\textbf{d}} & \uncover<4->{KPL-IFF: + \hl{limited signaling}, $C_N$ becomes
      non-signaling at rate $\varphi k_p$} & \uncover<4->{\yes{} all four}
  \end{tabular}
  \source{Lever et al., \emph{PNAS} (2016)}
\end{frame}

\begin{frame}{Is the mechanism unique?}
  \centering
  \regionfig[height=0.32\textheight]{58}{61}{120}{798}{324}
  \vskip1mm
  \begin{columns}[T]
    \begin{column}{0.24\textwidth}
      \centering
      \onslide<2->
      \regionfig[height=0.28\textheight]{58}{89}{347}{258}{529}
    \end{column}
    \begin{column}{0.26\textwidth}
      \centering
      \onslide<3->
      \regionfig[height=0.3\textheight]{59}{421}{333}{657}{542}
    \end{column}
    \begin{column}{0.18\textwidth}
      \centering
      \onslide<4->
      \regionfig[height=0.3\textheight]{60}{507}{308}{643}{538}
    \end{column}
    \begin{column}{0.28\textwidth}
      \onslide<1->
      \small
      \onslide<2->{Enumerate every network; for each, sample $10^6$ parameter sets and
      optimize the best 15.}
      \vskip1mm
      \onslide<3->{Compatible networks (1 of 304; 274 of 26,069) \hl{all have a feedforward}
      architecture.}
      \vskip1mm
      \onslide<4->{Feedback alone cannot produce the bell-shaped dose--response
      (shown analytically).}
    \end{column}
  \end{columns}
  \source{Lever et al., \emph{PNAS} (2016)}
\end{frame}

\begin{frame}{Possible relationship to known components}
  \centering
  \maskedfig[height=0.46\textheight]{62}{66}{133}{888}{523}{%
    \hide{1}{345}{130}{892}{526}%
    \hide{1-2}{630}{130}{892}{285}%
    \hide{1-3}{345}{285}{892}{526}}
  \vskip1mm
  \raggedright\footnotesize
  \begin{itemize}
    \setlength\itemsep{0.5mm}
    \item<2-> Kinetic proofreading: sequential ITAM phosphorylation and ZAP-70 recruitment.
    \item<3-> Limited signaling: TCR ubiquitination, or a membrane environment incompatible
          with signaling.
    \item<4-> Incoherent feedforward: LAT recruits SOS (activates Ras) and DOK1/2 with RasGAP
          (inactivates Ras).
  \end{itemize}
  \source{Lever et al., \emph{PNAS} (2016)}
\end{frame}

\begin{frame}{Key findings}
  \begin{itemize}
    \setlength\itemsep{2mm}
    \item<1-> Cytokine dose--response measured over a wide range of dose \emph{and}
          affinity.
    \item<2-> Systematic model rejection identifies a \hl{minimal model} that reproduces the
          observed behavior. It requires
          \begin{itemize}
            \item kinetic proofreading, for sensitivity to affinity
            \item limited signaling, which caps the response to long-lived complexes
            \item an incoherent feedforward loop, which gives the bell shape
          \end{itemize}
    \item<3-> The model predicted, and experiments confirmed, that the effect of
          co-presenting a weak ligand depends on dose.
  \end{itemize}
  \vskip3mm
  \onslide<4->
  \discuss{What did a minimal phenotypic model buy here that a detailed molecular model
  (like the 129-rule TCR model) could not? What did it give up?}
\end{frame}

\begin{frame}{The story continues}
  \small
  \begin{itemize}
    \setlength\itemsep{2.5mm}
    \item<1-> \textbf{Ganti et al.\ (2020)}: how the T cell signaling network processes
          information to discriminate between self and agonist ligands.
    \item<2-> \textbf{Pettmann et al.\ (2021)}: careful measurements of the discriminatory
          power of the TCR.
    \item<3-> \textbf{Britain, Town, and Weiner (2022)}: kinetic proofreading is
          progressively enhanced from receptor activation to DAG generation.
  \end{itemize}
  \vskip3mm
  \onslide<4->
  Each combines quantitative measurements with a model and asks which mechanism the data
  allow. That is the loop from the start of the lecture.
  \source{Ganti et al., \emph{PNAS} \textbf{117}, 26020 (2020); Pettmann et al., \emph{eLife} \textbf{10}, e67092 (2021); Britain et al., \emph{eLife} \textbf{11}, e75263 (2022)}
\end{frame}

% ======================================================================
% Part 6  (~12 min)
\section{How far can mechanistic models scale?}

\begin{frame}{A pan-cancer mechanistic model of drug response}
  \begin{columns}[c]
    \begin{column}{0.46\textwidth}
      \centering
      \regionfig[height=0.7\textheight]{68}{48}{112}{416}{540}
    \end{column}
    \begin{column}{0.5\textwidth}
      \small
      \begin{itemize}
        \setlength\itemsep{1.5mm}
        \item<1-> ERBB, RAS, PI3K--AKT, and MYC/AP1 pathways, assembled from
              ConsensusPathDB
        \item<2-> 108 genes, 36 activating mutations (7 of the 10 most frequent cancer
              drivers)
        \item<3-> 1,228 species in 4 compartments, 2,686 reactions
        \item<4-> 7 kinase inhibitors, 4.3 targets per drug on average
        \item<5-> Phenotype: cell viability modeled as a weighted sum of pathway activities
      \end{itemize}
    \end{column}
  \end{columns}
  \source{Fr\"ohlich et al., \emph{Cell Systems} \textbf{7}, 567 (2018)}
\end{frame}

\begin{frame}{Larger than previous models by an order of magnitude}
  \begin{columns}[c]
    \begin{column}{0.56\textwidth}
      \centering
      \regionfig[height=0.76\textheight]{71}{237}{112}{723}{511}
    \end{column}
    \begin{column}{0.4\textwidth}
      \small
      \begin{itemize}
        \setlength\itemsep{2mm}
        \item<1-> Compared with curated BioModels and earlier signaling models
              (Kholodenko 1999, Chen 2009, \ldots).
        \item<2-> Most species are molecular \emph{complexes}: combinatorial complexity
              again.
      \end{itemize}
    \end{column}
  \end{columns}
  \source{Fr\"ohlich et al., \emph{Cell Systems} \textbf{7}, 567 (2018)}
\end{frame}

\begin{frame}{Personalizing the model with genomic data}
  \begin{columns}[c]
    \begin{column}{0.48\textwidth}
      \small
      \begin{itemize}
        \setlength\itemsep{2mm}
        \item<1-> 120 cancer cell lines from 5 tissues.
        \item<2-> Gene expression sets the synthesis rate of each protein and functional
              variant in each cell line.
        \item<3-> All other parameters are shared across cell lines, so every cell line
              adds training data for the same parameters.
      \end{itemize}
    \end{column}
    \begin{column}{0.48\textwidth}
      \centering
      \regionfig[height=0.78\textheight]{72}{560}{109}{912}{501}
    \end{column}
  \end{columns}
  \source{Fr\"ohlich et al., \emph{Cell Systems} \textbf{7}, 567 (2018)}
\end{frame}

\begin{frame}{Large-scale parameter estimation}
  \begin{columns}[c]
    \begin{column}{0.4\textwidth}
      \centering
      \regionfig[width=\linewidth]{73}{10}{162}{319}{418}
    \end{column}
    \begin{column}{0.56\textwidth}
      \small
      \begin{itemize}
        \setlength\itemsep{1.5mm}
        \item<1-> $>$4,100 unknown parameters
        \item<2-> Trained on viability of 120 cell lines $\times$ 7 drugs $\times$ up to 9
              concentrations: 6,900 conditions, 5-fold cross-validation
        \item<3-> Problem: one gradient of the objective would take $>$6 CPU years
        \item<4-> Solution: $>$5 orders of magnitude speed-up from
              \begin{itemize}
                \item sparse linear solvers for ODE integration
                \item adjoint sensitivity analysis
                \item coarse-grained parallelization
              \end{itemize}
      \end{itemize}
    \end{column}
  \end{columns}
  \source{Fr\"ohlich et al., \emph{Cell Systems} \textbf{7}, 567 (2018)}
\end{frame}

\begin{frame}{Fits and predictions}
  \centering
  \only<1|handout:1>{\regionfig[height=0.62\textheight]{75}{108}{136}{852}{470}}%
  \only<2|handout:2>{\regionfig[height=0.56\textheight]{77}{96}{147}{846}{438}}%
  \only<3|handout:3>{\regionfig[height=0.62\textheight]{80}{177}{112}{783}{463}}%
  \vskip1mm
  \raggedright\small
  \only<1|handout:1>{Quantitative agreement across cell lines, drugs, and concentrations.}%
  \only<2|handout:2>{Prediction accuracy on held-out data is comparable to established
    machine-learning (statistical) methods.}%
  \only<3|handout:3>{The model, trained on single drugs, predicts responses to drug
    combinations.}
  \source{Fr\"ohlich et al., \emph{Cell Systems} \textbf{7}, 567 (2018)}
\end{frame}

% TODO (revisit): check whether the frontier in scale has moved since 2018.
\begin{frame}{Since 2018: has the frontier moved?}
  \begin{columns}[T]
    \begin{column}{0.55\textwidth}
      \footnotesize
      \begin{itemize}
        \setlength\itemsep{1.2mm}
        \item<1-> \textbf{Scale}: calibrated models are still about $10^3$ species and
              $10^3$--$10^4$ parameters. Example: SPARCED, a single-cell pan-cancer
              model of proliferation and death signaling (Erdem et al., 2022).
        \item<2-> \textbf{Infrastructure}: standard formats and tools make
              parameter-estimation problems reusable, e.g.\ PEtab (Schmiester et al., 2021).
        \item<3-> \textbf{The bottleneck} is building and curating models: reading the
              literature, choosing mechanisms, and assembling rules.
        \item<4-> \textbf{LLMs} may change exactly that step: drafting rules from papers,
              proposing alternative mechanisms, checking models against data.
      \end{itemize}
    \end{column}
    \begin{column}{0.41\textwidth}
      \vskip2mm
      \onslide<5->
      \discuss{If constructing a model became cheap, what would limit mechanistic
      modeling: data, parameter identifiability, or validation? How would you check a
      model that an LLM wrote?}
    \end{column}
  \end{columns}
  \source{Erdem et al., \emph{Nat.\ Commun.}\ \textbf{13}, 3555 (2022); Schmiester et al., \emph{PLoS Comput.\ Biol.}\ \textbf{17}, e1008646 (2021)}
\end{frame}

% ======================================================================
\begin{frame}{Summary}
  \begin{itemize}[<+->]
    \setlength\itemsep{2mm}
    \item Mechanistic models turn assumptions about mechanism into testable,
          quantitative predictions, and close the loop with experiment.
    \item \textbf{Simple models} can be analyzed rigorously and expose core principles:
          kinetic proofreading, feedback, feedforward.
    \item \textbf{Detailed models} map known molecular components onto mechanisms.
          Rule-based modeling handles the combinatorial complexity.
    \item \textbf{Calibration} to data is the main bottleneck. New algorithms now fit
          models with thousands of parameters to thousands of conditions.
    \item Model selection against comprehensive data can identify the \textbf{minimal
          architecture} a system must have.
  \end{itemize}
\end{frame}

\begin{frame}{References}
  \tiny
  \setlength\parskip{0.6mm}
  \begin{columns}[T]
    \begin{column}{0.49\textwidth}
      H. Kitano, ``Systems biology: a brief overview,'' \emph{Science} 295, 1662 (2002).\par
      Y. Lazebnik, ``Can a biologist fix a radio?,'' \emph{Cancer Cell} 2, 179 (2002).\par
      T. W. McKeithan, ``Kinetic proofreading in T-cell receptor signal transduction,''
        \emph{PNAS} 92, 5042 (1995).\par
      B. N. Kholodenko et al., ``Quantification of short term signaling by the epidermal
        growth factor receptor,'' \emph{J.\ Biol.\ Chem.}\ 274, 30169 (1999).\par
      A. Hoffmann et al., ``The I$\kappa$B--NF-$\kappa$B signaling module: temporal control
        and selective gene activation,'' \emph{Science} 298, 1241 (2002).\par
      J. R. Faeder et al., ``Investigation of early events in Fc$\varepsilon$RI-mediated
        signaling using a detailed mathematical model,'' \emph{J.\ Immunol.}\ 170, 3769 (2003).\par
      B. Goldstein, J. R. Faeder, and W. S. Hlavacek, ``Mathematical and computational models
        of immune-receptor signalling,'' \emph{Nat.\ Rev.\ Immunol.}\ 4, 445 (2004).\par
      M. L. Blinov et al., ``BioNetGen: software for rule-based modeling of signal
        transduction,'' \emph{Bioinformatics} 20, 3289 (2004).\par
      G. Altan-Bonnet and R. N. Germain, ``Modeling T cell antigen discrimination based on
        feedback control of digital ERK responses,'' \emph{PLoS Biol.}\ 3, e356 (2005).\par
      M. W. Sneddon, J. R. Faeder, and T. Emonet, ``Efficient modeling, simulation and
        coarse-graining of biological complexity with NFsim,'' \emph{Nat.\ Methods} 8, 177 (2011).\par
      J. A. P. Sekar and J. R. Faeder, ``Rule-based modeling of signal transduction: a
        primer,'' \emph{Methods Mol.\ Biol.}\ 880, 139 (2012).\par
      N. Miskov-Zivanov et al., ``The duration of T cell stimulation is a critical determinant
        of cell fate and plasticity,'' \emph{Sci.\ Signal.}\ 6, ra97 (2013).\par
      L. A. Chylek et al., ``Phosphorylation site dynamics of early T-cell receptor
        signaling,'' \emph{PLOS One} 9, e104240 (2014).\par
    \end{column}
    \begin{column}{0.49\textwidth}
      M. Lever et al., ``Phenotypic models of T cell activation,'' \emph{Nat.\ Rev.\ Immunol.}\
        14, 619 (2014).\par
      W. F. Hawse et al., ``Differential regulation of PTEN by TCR, Akt, and FoxO1 controls
        CD4$^+$ T cell fate decisions,'' \emph{J.\ Immunol.}\ 194, 4615 (2015).\par
      L. A. Harris et al., ``BioNetGen 2.2: advances in rule-based modeling,''
        \emph{Bioinformatics} 32, 3366 (2016).\par
      M. Lever et al., ``Architecture of a minimal signaling pathway explains the T-cell
        response to a 1 million-fold variation in antigen affinity and dose,''
        \emph{PNAS} 113, E6630 (2016).\par
      F. Fr\"ohlich et al., ``Efficient parameter estimation enables the prediction of drug
        response using a mechanistic pan-cancer pathway model,'' \emph{Cell Syst.}\ 7, 567 (2018).\par
      J. A. P. Sekar and J. R. Faeder, ``An introduction to rule-based modeling of immune
        receptor signaling,'' in \emph{Systems Immunology}, CRC Press (2019).\par
      R. S. Ganti et al., ``How the T cell signaling network processes information to
        discriminate between self and agonist ligands,'' \emph{PNAS} 117, 26020 (2020).\par
      L. Schmiester et al., ``PEtab: interoperable specification of parameter estimation
        problems in systems biology,'' \emph{PLoS Comput.\ Biol.}\ 17, e1008646 (2021).\par
      J. Pettmann et al., ``The discriminatory power of the T cell receptor,''
        \emph{eLife} 10, e67092 (2021).\par
      D. M. Britain, J. P. Town, and O. D. Weiner, ``Progressive enhancement of kinetic
        proofreading in T cell antigen discrimination from receptor activation to DAG
        generation,'' \emph{eLife} 11, e75263 (2022).\par
      C. Erdem et al., ``A scalable, open-source implementation of a large-scale mechanistic
        model for single cell proliferation and death signaling,''
        \emph{Nat.\ Commun.}\ 13, 3555 (2022).\par
    \end{column}
  \end{columns}
\end{frame}

\end{document}
